\section{教学评价体系}

本案例采用过程性评价（40\%）、成果性评价（40\%）与反思性评价（20\%）相结合的多元化评价体系，如表1所示。

\begin{table}[htbp]
\centering
\caption{教学评价体系}
\begin{tabular}{p{3.5cm}p{2cm}p{8cm}}
\toprule
\textbf{评价维度} & \textbf{权重} & \textbf{评价要点} \\
\midrule
AI交互日志 & 20\% & 提问的精准性（问题清晰具体）、逻辑性（问题链条完整）、批判性（识别AI不足并修正） \\
\midrule
代码调试记录 & 20\% & 代码理解（准确解释功能原理）、错误修正（独立发现修正错误）、创新改进（优化适应具体情况） \\
\midrule
选址方案报告 & 30\% & 科学性（方法正确、数据规范、权重有据）、创新性（独特视角或方法改进）、图件规范性（要素完整、图例清晰） \\
\midrule
中学教学案例设计 & 10\% & 学科契合度（紧扣课标、体现核心素养）、教学可行性（难度适中符合认知）、创新性（独特设计思路） \\
\midrule
反思报告 & 20\% & 反思深度（深入分析AI优势局限）、实践关联（与具体案例结合）、未来展望（建设性建议） \\
\bottomrule
\end{tabular}
\end{table}

评价实施要点：（1）AI交互日志采用档案袋评价，学生全程记录，教师定期检查反馈；（2）代码调试记录重点考查理解深度而非运行成功率；（3）选址方案报告需附GIS操作截图与权重确定论证；\（4）中学教学案例设计强调与高中课标的对接；（5）反思报告采用答辩形式，学生需现场回答教师提问。